\scanpage{226}
% continuation; the bracketed draft is crossed through with a wavy stroke
dernière, on \deleted{[peut se ramener, en décomposant $Z_{s_0}$ en \deleted{\unsure{petits arcs}} \inserted{petits} chemins $a_i^0\longrightarrow a^0_{i+1}$, et transportant \deleted{par} horizontalement les $a_i^0$, \deleted{\unsure{à}}]} va prouver la relation plus générale
\[
 D_s\int_{a(s)}^{b(s)}\omega_s=\int_{a(s)}^{b(s)}(\overline D\omega)_s,
 \tag*{(13)}
\]
valable pour des \unsure{limites|éléments|chemins} d'intégration $a(s)$, $b(s)$ qui sont des \uline{sections} \uline{horizontales} pour la connexion. [N.B. Bien \deleted{entendu}, on suppose, pour donner un sens à cette formule, qu'on choisit les $\overline{\mathrm m}^{\mathrm s}$ $C_s$ au deux mbs, donc (13) s'écrit avec tte précision
\[
 D_s\int_{C_s}\omega_s=\int_{C_s}(\overline D\omega)_s\qquad].
 \tag*{(13 bis)}
\]

% subdivision, ending in the middle of a sentence
Pour prouver (13), on est ramené, en \inserted{se plaçant au voisinage de $s_0\in S$ et} décomposant le \unsure{domaine|chemin} $\int_{a(s)}^{b(s)}$, au cas où $a(s_0),b(s_0)$ sont assez voisins, de sorte qu'on peut se placer dans un ouvert $U$ où l'on a $\omega=d\varphi$. Mais alors le $1^{\mathrm{er}}$ mb est $D_s\bigl(\varphi(b(s))-\varphi(a(s))\bigr)$, le
% archive sheet number: 225
